% ============================================================
%  EV Charge DZ - Shared Database API Report
%  Compile:  pdflatex Shared_API_Report.tex  (run twice)
% ============================================================
\documentclass[12pt,a4paper]{report}

% -- Encoding & fonts ----------------------------------------
\usepackage[T1]{fontenc}
\usepackage[utf8]{inputenc}
\usepackage{lmodern}
\usepackage{microtype}
\usepackage{amssymb}

% -- Page geometry -------------------------------------------
\usepackage[top=2.5cm, bottom=2.8cm, left=2.8cm, right=2.8cm]{geometry}

% -- Colours -------------------------------------------------
\usepackage{xcolor}
\definecolor{brandblue}{HTML}{1E3A8A}
\definecolor{lightblue}{HTML}{EFF6FF}
\definecolor{midblue}{HTML}{BFDBFE}
\definecolor{steelblue}{HTML}{2563EB}
\definecolor{codebg}{HTML}{F1F5F9}
\definecolor{greenbg}{HTML}{F0FDF4}
\definecolor{greenfg}{HTML}{14532D}
\definecolor{labelgray}{HTML}{475569}
\definecolor{rowalt}{HTML}{F8FAFC}
\definecolor{tablehead}{HTML}{1E3A8A}
\definecolor{orangebg}{HTML}{FFF7ED}
\definecolor{orangefg}{HTML}{C2410C}

% -- tcolorbox -----------------------------------------------
\usepackage[most]{tcolorbox}
\tcbuselibrary{listings, breakable, skins}

% -- listings ------------------------------------------------
\usepackage{listings}
\lstdefinestyle{codestyle}{
    basicstyle    = \ttfamily\footnotesize,
    breaklines    = true,
    keepspaces    = true,
    showstringspaces = false,
    tabsize       = 2,
    upquote       = true,
}
\lstset{style=codestyle}

\newtcblisting{codebox}{
    colback      = codebg,
    colframe     = midblue,
    listing only,
    listing options = {style=codestyle, backgroundcolor=\color{codebg}},
    arc          = 3pt,
    left         = 6pt, right = 6pt, top = 4pt, bottom = 4pt,
    breakable,
}

\newtcblisting{greenbox}{
    colback      = greenbg,
    colframe     = greenfg!40,
    listing only,
    listing options = {style=codestyle, backgroundcolor=\color{greenbg}},
    arc          = 3pt,
    left         = 6pt, right = 6pt, top = 4pt, bottom = 4pt,
    breakable,
}

% -- Tables --------------------------------------------------
\usepackage{booktabs}
\usepackage{tabularx}
\usepackage{longtable}
\usepackage{colortbl}
\usepackage{array}

% -- Lists ---------------------------------------------------
\usepackage{enumitem}
\setlist[itemize]{leftmargin=1.4em, itemsep=2pt, topsep=4pt, parsep=0pt}
\setlist[enumerate]{leftmargin=1.6em, itemsep=2pt, topsep=4pt, parsep=0pt}

% -- Section styling -----------------------------------------
\usepackage{titlesec}
\titleformat{\chapter}[display]
  {\normalfont\Large\bfseries\color{brandblue}}
  {\color{brandblue!60}\MakeUppercase{\chaptertitlename}\ \thechapter}
  {8pt}{\LARGE}
  [\vspace{-4pt}\color{midblue}\rule{\linewidth}{1pt}]
\titlespacing*{\chapter}{0pt}{0pt}{14pt}

\titleformat{\section}
  {\normalfont\large\bfseries\color{brandblue}}
  {\thesection}{0.8em}{}
  [\vspace{-4pt}\color{midblue}\rule{\linewidth}{0.5pt}]
\titlespacing*{\section}{0pt}{16pt}{6pt}

\titleformat{\subsection}
  {\normalfont\normalsize\bfseries\color{steelblue}}
  {\thesubsection}{0.8em}{}
\titlespacing*{\subsection}{0pt}{10pt}{4pt}

% -- Header / footer -----------------------------------------
\usepackage{fancyhdr}
\pagestyle{fancy}
\fancyhf{}
\setlength{\headheight}{14pt}
\fancyhead[L]{\small\color{brandblue}\textbf{EV Charge DZ}\enspace\textcolor{midblue}{|}\enspace Shared Database API Report}
\fancyhead[R]{\small\color{labelgray}May 2026}
\fancyfoot[C]{\small\color{labelgray}--- \thepage\ ---}
\renewcommand{\headrulewidth}{0.5pt}
\renewcommand{\headrule}{\color{midblue}\hrule width\headwidth height\headrulewidth\vskip-\headrulewidth}
\renewcommand{\footrulewidth}{0pt}

% -- Spacing -------------------------------------------------
\usepackage{parskip}
\setlength{\parindent}{0pt}
\setlength{\parskip}{5pt}

% -- Hyperref -----------------------------------------------
\usepackage[
  colorlinks  = true,
  linkcolor   = brandblue,
  urlcolor    = steelblue,
  pdftitle    = {EV Charge DZ - Shared Database API Report},
  pdfauthor   = {Fodhil \& Chibani},
  pdfpagelabels = false,
]{hyperref}

% -- Convenience macros --------------------------------------
\newcommand{\file}[1]{\texttt{\small #1}}
\newcommand{\done}{\textcolor{greenfg}{$\checkmark$}}
\newcommand{\endpoint}[1]{\texttt{\textbf{\small #1}}}

% ============================================================
\begin{document}
\pagenumbering{Alph}

% -- TITLE PAGE ---------------------------------------------
\begin{titlepage}
  \centering
  \vspace*{2.5cm}

  {\fontsize{42}{50}\selectfont\bfseries\color{brandblue} EV Charge DZ}

  \vspace{0.5cm}
  {\LARGE Shared Database API Report}

  \vspace{0.3cm}
  \textcolor{midblue}{\rule{0.6\linewidth}{1pt}}

  \vspace{0.5cm}
  {\large\color{labelgray} Web Platform $\boldsymbol{\leftrightarrow}$ Shared MySQL Database $\boldsymbol{\leftrightarrow}$ Kotlin App}

  \vspace{2cm}

  \begin{tcolorbox}[
    colback  = lightblue,
    colframe = midblue,
    arc      = 6pt,
    width    = 0.75\linewidth,
    halign   = center,
  ]
    \begin{tabular}{ll}
      \textbf{\color{brandblue}Project}   & EV Charge DZ \\[3pt]
      \textbf{\color{brandblue}Authors}   & Fodhil Yacine \& Chibani \\[3pt]
      \textbf{\color{brandblue}Date}      & May 2026 \\[3pt]
      \textbf{\color{brandblue}Backend}   & Node.js / Express (TypeScript) \\[3pt]
      \textbf{\color{brandblue}Database}  & MySQL 8.0 — \texttt{ev\_charge\_dz} \\[3pt]
      \textbf{\color{brandblue}API Port}  & \texttt{3001} \\
    \end{tabular}
  \end{tcolorbox}

  \vfill
  {\small\color{labelgray} This document summarises the API changes made to support
  the shared-database architecture between the admin web platform and the Kotlin
  mobile application.}
\end{titlepage}

% -- TOC ----------------------------------------------------
\pagenumbering{roman}
\tableofcontents
\clearpage
\pagenumbering{arabic}

% ============================================================
\chapter{Context and Architecture}
% ============================================================

\section{Goal}

The project uses a \textbf{single shared MySQL database} (\texttt{ev\_charge\_dz}) accessed by two independent clients:

\begin{itemize}
  \item \textbf{Web Platform} — React admin dashboard managed by Fodhil Yacine.
  \item \textbf{Kotlin Mobile App} — Android application managed by Chibani.
\end{itemize}

Both clients communicate with the same Express backend running on port \texttt{3001}.
The teacher's requirement was that every path must return a \textbf{JSON response}, never HTML,
so that the Kotlin app can consume the same API as the web platform.

\section{Before the Change}

All routes were mounted directly at the root:

\begin{codebox}
app.use('/auth',     authRoutes);
app.use('/stations', stationRoutes);
app.use('/sessions', sessionRoutes);
// ... etc.
\end{codebox}

There was no clear contract between the web platform and the mobile app.
The Kotlin app had no dedicated authentication or data endpoints.

\section{After the Change}

\begin{tcolorbox}[colback=lightblue, colframe=midblue, arc=4pt]
\textbf{All routes now live under \texttt{/api/}.}
The web platform uses \texttt{/api/*} and the Kotlin app uses \texttt{/api/app/*}.
Both return pure JSON. No HTML is ever served by the backend.
\end{tcolorbox}

\begin{center}
\begin{tcolorbox}[colback=codebg, colframe=midblue, arc=4pt, width=0.9\linewidth]
\begin{lstlisting}
http://localhost:3001
  /health                   -- health check (no auth)
  /api/auth/*               -- web platform authentication
  /api/stations/*           -- web platform station management
  /api/sessions/*           -- web platform session listing
  /api/tickets/*            -- maintenance tickets
  /api/billing/*            -- billing records
  /api/alerts/*             -- system alerts
  /api/dashboard/*          -- KPI data
  /api/users/*              -- dashboard user management
  /api/tenants/*            -- tenant management
  /api/connectors/*         -- connector management
  /api/app/auth/*           -- Kotlin app authentication
  /api/app/stations/*       -- Kotlin app station map
  /api/app/sessions/*       -- Kotlin app session history
\end{lstlisting}
\end{tcolorbox}
\end{center}

% ============================================================
\chapter{Changes Made — Web Platform Side (Yacine)}
% ============================================================

\section{Files Modified}

\subsection{\file{server/src/index.ts} — Route prefix}

Added \texttt{/api} prefix to all ten existing route handlers and imported the three
new app route files.

\begin{greenbox}
// Web platform routes
app.use('/api/auth',       authRoutes);
app.use('/api/tenants',    tenantRoutes);
app.use('/api/users',      userRoutes);
app.use('/api/stations',   stationRoutes);
app.use('/api/connectors', connectorRoutes);
app.use('/api/sessions',   sessionRoutes);
app.use('/api/tickets',    ticketRoutes);
app.use('/api/billing',    billingRoutes);
app.use('/api/alerts',     alertRoutes);
app.use('/api/dashboard',  dashboardRoutes);

// Kotlin app routes
app.use('/api/app/auth',     appAuthRoutes);
app.use('/api/app/stations', appStationsRoutes);
app.use('/api/app/sessions', appSessionsRoutes);
\end{greenbox}

\subsection{\file{server/src/middleware/auth.ts} — App user middleware}

Added a second JWT payload interface (\texttt{AppAuthPayload}) and a new
\texttt{requireAppAuth} middleware for Kotlin app users.
Dashboard users and app users have separate identities in the database
(\texttt{users} table vs \texttt{app\_users} table).

\begin{greenbox}
export interface AppAuthPayload {
  appUserId: string;
  email:     string;
}

export function requireAppAuth(req, res, next): void {
  // Verifies Bearer token, populates req.appUser
}
\end{greenbox}

\subsection{\file{project/.env} — React base URL}

Updated \texttt{VITE\_API\_BASE\_URL} so the React dashboard continues to work
without changing any path in the frontend code:

\begin{greenbox}
# Before
VITE_API_BASE_URL=http://localhost:3001

# After
VITE_API_BASE_URL=http://localhost:3001/api
\end{greenbox}

\section{New Files Created}

\subsection{\file{server/src/routes/appAuth.ts}}

Handles \texttt{app\_users} registration and login (completely separate from
dashboard staff accounts).

\begin{center}
\begin{tabularx}{0.9\linewidth}{>{\ttfamily\small}l l X}
  \toprule
  \rowcolor{tablehead}\color{white}\textbf{Method} & \color{white}\textbf{Path} & \color{white}\textbf{Description} \\
  \midrule
  POST & \texttt{/api/app/auth/register} & Create new app user account \\
  \rowcolor{rowalt}
  POST & \texttt{/api/app/auth/login}    & Login, returns JWT token \\
  \bottomrule
\end{tabularx}
\end{center}

\subsection{\file{server/src/routes/appStations.ts}}

Returns station data shaped for the Kotlin map view using the existing
\texttt{vw\_app\_station\_map} database view. \textbf{No authentication required}
so the map loads before the user logs in.

\begin{center}
\begin{tabularx}{0.9\linewidth}{>{\ttfamily\small}l l X}
  \toprule
  \rowcolor{tablehead}\color{white}\textbf{Method} & \color{white}\textbf{Path} & \color{white}\textbf{Description} \\
  \midrule
  GET & \texttt{/api/app/stations}      & All available stations (public) \\
  \rowcolor{rowalt}
  GET & \texttt{/api/app/stations/\{id\}} & Single station detail (public) \\
  \bottomrule
\end{tabularx}
\end{center}

\subsection{\file{server/src/routes/appSessions.ts}}

Returns the authenticated user's personal charge history using the existing
\texttt{vw\_app\_session\_history} database view.

\begin{center}
\begin{tabularx}{0.9\linewidth}{>{\ttfamily\small}l l X}
  \toprule
  \rowcolor{tablehead}\color{white}\textbf{Method} & \color{white}\textbf{Path} & \color{white}\textbf{Description} \\
  \midrule
  GET & \texttt{/api/app/sessions} & Logged-in user's session history \\
  \bottomrule
\end{tabularx}
\end{center}

Supports optional query parameters: \texttt{limit} (default 50, max 200) and \texttt{offset} (default 0).

% ============================================================
\chapter{For the Kotlin App Developer (Chibani)}
% ============================================================

\section{Base URL}

Replace \texttt{<server-ip>} with the IP address of the machine running the backend
(use \texttt{localhost} or \texttt{10.0.2.2} for the Android emulator, or the actual
LAN IP for a physical device).

\begin{codebox}
val BASE_URL = "http://<server-ip>:3001/api/app"
\end{codebox}

\begin{tcolorbox}[colback=orangebg, colframe=orangefg!40, arc=4pt]
  \textbf{Android emulator note:} \texttt{10.0.2.2} maps to \texttt{localhost} on the host machine.
  Physical device on the same Wi-Fi: use the host machine's local IP (e.g.\ \texttt{192.168.x.x}).
\end{tcolorbox}

\section{Authentication Flow}

\subsection{Register}

\begin{codebox}
POST http://<server-ip>:3001/api/app/auth/register
Content-Type: application/json

{
  "email":     "user@example.com",
  "password":  "SecurePass123",
  "full_name": "Ahmed Benali",
  "phone":     "+213555000000"   // optional
}
\end{codebox}

\textbf{Response 201:}
\begin{codebox}
{
  "token": "<jwt-token>",
  "user": {
    "appUserId": "uuid",
    "email":     "user@example.com",
    "full_name": "Ahmed Benali",
    "phone":     "+213555000000"
  }
}
\end{codebox}

\subsection{Login}

\begin{codebox}
POST http://<server-ip>:3001/api/app/auth/login
Content-Type: application/json

{
  "email":    "user@example.com",
  "password": "SecurePass123"
}
\end{codebox}

\textbf{Response 200:}
\begin{codebox}
{
  "token": "<jwt-token>",
  "user": {
    "appUserId": "uuid",
    "email":     "user@example.com",
    "full_name": "Ahmed Benali",
    "phone":     "+213555000000"
  }
}
\end{codebox}

\textbf{Save the token} in SharedPreferences or DataStore. Every authenticated
request must include it in the header:

\begin{codebox}
Authorization: Bearer <jwt-token>
\end{codebox}

\section{Station Map}

\subsection{Get All Stations (public — no token needed)}

\begin{codebox}
GET http://<server-ip>:3001/api/app/stations
\end{codebox}

\textbf{Response 200} — array of station objects:
\begin{codebox}
[
  {
    "id":                   "uuid",
    "name":                 "Station Alger Centre",
    "address":              "Rue Didouche Mourad",
    "city":                 "Alger",
    "lat":                  36.7525,
    "lng":                  3.0420,
    "status":               "available",
    "available":            1,
    "max_power_kw":         50,
    "power":                "50 kW",
    "price":                "25 DZD/kWh",
    "total_connectors":     4,
    "available_connectors": 2,
    "connector_types":      "CCS, Type2"
  },
  ...
]
\end{codebox}

\subsection{Get Single Station}

\begin{codebox}
GET http://<server-ip>:3001/api/app/stations/{id}
\end{codebox}

Returns a single station object with the same fields as above, or \texttt{404} if not found.

\section{Session History (requires token)}

\begin{codebox}
GET http://<server-ip>:3001/api/app/sessions
Authorization: Bearer <jwt-token>
\end{codebox}

Optional query parameters:

\begin{center}
\begin{tabularx}{0.75\linewidth}{l l X}
  \toprule
  \rowcolor{tablehead}\color{white}\textbf{Parameter} & \color{white}\textbf{Default} & \color{white}\textbf{Description} \\
  \midrule
  \texttt{limit}  & 50  & Number of records to return (max 200) \\
  \rowcolor{rowalt}
  \texttt{offset} & 0   & Skip this many records (for pagination) \\
  \bottomrule
\end{tabularx}
\end{center}

\textbf{Response 200} — array of session objects:
\begin{codebox}
[
  {
    "session_id":      "uuid",
    "app_user_id":     "uuid",
    "station_id":      "uuid",
    "station_name":    "Station Alger Centre",
    "station_address": "Rue Didouche Mourad",
    "connector_id":    "uuid",
    "connector_type":  "CCS",
    "max_power_kw":    50,
    "payment_method":  "card",
    "start_time":      "2026-05-01T10:30:00Z",
    "end_time":        "2026-05-01T11:15:00Z",
    "duration_min":    45,
    "energy_kwh":      22.5,
    "cost_dzd":        562.5,
    "status":          "completed"
  },
  ...
]
\end{codebox}

\section{Error Responses}

All endpoints return errors in a consistent format:

\begin{codebox}
{ "error": "descriptive message" }
\end{codebox}

\begin{center}
\begin{tabularx}{0.85\linewidth}{l l X}
  \toprule
  \rowcolor{tablehead}\color{white}\textbf{HTTP Code} & \color{white}\textbf{Meaning} & \color{white}\textbf{When it happens} \\
  \midrule
  400 & Bad Request     & Missing required fields \\
  \rowcolor{rowalt}
  401 & Unauthorized    & No token or invalid token \\
  403 & Forbidden       & Account suspended \\
  \rowcolor{rowalt}
  404 & Not Found       & Station ID does not exist \\
  409 & Conflict        & Email already registered \\
  \rowcolor{rowalt}
  500 & Server Error    & Unexpected backend error \\
  \bottomrule
\end{tabularx}
\end{center}

\section{Recommended Kotlin Integration Steps}

\begin{enumerate}
  \item Add \textbf{Retrofit} + \textbf{OkHttp} + \textbf{Gson} (or Moshi) to \file{build.gradle}.
  \item Create a \textbf{Retrofit service interface} with the endpoints above.
  \item Add an \textbf{OkHttp interceptor} that reads the JWT from DataStore and injects
        the \texttt{Authorization} header automatically on every request.
  \item Create \textbf{data classes} matching the JSON fields (use \texttt{@SerializedName} for
        snake\_case fields like \texttt{station\_name}, \texttt{start\_time}, etc.).
  \item Map \file{vw\_app\_station\_map} fields to a \texttt{ChargingStation} data class and
        \file{vw\_app\_session\_history} fields to a \texttt{SessionRecord} data class.
  \item For the map: load stations on \texttt{onMapReady} — no login required.
  \item For session history: check token validity first; redirect to login if 401.
\end{enumerate}

\section{Quick Reference Card}

\begin{center}
\begin{tabularx}{0.95\linewidth}{l >{\ttfamily\small}l l l}
  \toprule
  \rowcolor{tablehead}
  \color{white}\textbf{Feature} &
  \color{white}\textbf{Endpoint} &
  \color{white}\textbf{Method} &
  \color{white}\textbf{Auth} \\
  \midrule
  Register          & /api/app/auth/register    & POST & None \\
  \rowcolor{rowalt}
  Login             & /api/app/auth/login       & POST & None \\
  Station map       & /api/app/stations         & GET  & None \\
  \rowcolor{rowalt}
  Station detail    & /api/app/stations/\{id\}  & GET  & None \\
  Session history   & /api/app/sessions         & GET  & Bearer token \\
  \bottomrule
\end{tabularx}
\end{center}

% ============================================================
\chapter{Summary}
% ============================================================

\begin{tcolorbox}[colback=greenbg, colframe=greenfg!40, arc=4pt]
\textbf{What was completed on the web platform side:}
\begin{itemize}
  \item[\done] All existing routes prefixed with \texttt{/api/} — web platform still works via updated \texttt{VITE\_API\_BASE\_URL}.
  \item[\done] App user JWT middleware added (\texttt{requireAppAuth}).
  \item[\done] \file{appAuth.ts}: register \& login for \texttt{app\_users}.
  \item[\done] \file{appStations.ts}: public station map endpoint using \texttt{vw\_app\_station\_map}.
  \item[\done] \file{appSessions.ts}: authenticated session history using \texttt{vw\_app\_session\_history}.
\end{itemize}
\end{tcolorbox}

\vspace{0.5cm}

\begin{tcolorbox}[colback=lightblue, colframe=midblue, arc=4pt]
\textbf{What Chibani needs to do on the Kotlin side:}
\begin{itemize}
  \item Configure Retrofit with base URL \texttt{http://<server-ip>:3001/api/app}.
  \item Implement register / login screens calling \texttt{/auth/register} and \texttt{/auth/login}.
  \item Save the returned JWT token in DataStore / SharedPreferences.
  \item Inject the token via an OkHttp interceptor on authenticated requests.
  \item Map \texttt{vw\_app\_station\_map} columns to a \texttt{ChargingStation} data class.
  \item Map \texttt{vw\_app\_session\_history} columns to a \texttt{SessionRecord} data class.
  \item Load stations (no auth) for the map; load sessions (with auth) for the history screen.
\end{itemize}
\end{tcolorbox}

\end{document}
